% ============================================================
%  第2章  系统总体方案设计
% ============================================================
\chapter{系统总体方案设计}

本章根据脑卒中后视觉障碍康复的临床需求，对头戴式眼部康复装置的功能需求、总体架构方案和关键技术指标进行了系统的分析，为后续的硬件选型与电路设计、软件架构和算法实现提供宏观层面的参考。

\section{功能需求分析}

\subsection{临床康复需求}

经由临床研究表明，脑卒中后视觉障碍的康复训练应满足以下几项基本原则~\cite{strokerehab2011}：训练任务应当具有明确的功能意义，可以通过具体的效验指标来评价疗效的好坏，不能只依靠主观感觉去判定效果是否理想，而且必须要给予病人即时的反馈信息（视觉、听觉、触觉等），这样才能够及时改正错误的行为并增进正确的动作方式，充分发挥患者的神经系统可塑性潜力并使其发展速度与之保持一致，不断地加大训练的强度使得患者神经可塑性的成长速度同它成长的速度相符合，力求使患者在日常生活中碰到的各种状况都能够包含进去，以提高康复效果的生态效度和日常生活迁移性。

根据以上原则，系统的训练功能应覆盖以下五种核心视觉能力。

注视稳定就是指把视线固定在一个地方。从解剖学的角度来说，它依靠前庭眼反射(VOR)的完好无损，并且由脑干和小脑进行细致地调节。脑卒中之后，由于中枢神经通路受到了损害，患者的视觉会出现无法保持固视的情况，在没有自发扫视行为的时候，眼球也会有轻微的移动的现象发生，这就造成了视觉信息的输入被中断了。注视稳定性是其他所有高级视觉功能的基础，只有眼睛一直盯着一个目标不动的时候，才能去做其他的诸如扫视寻找、平滑追随着或者空间注意这些事情。

快速扫视就是眼球在视野内从一个目标迅速转向另一个目标的能力。阅读、驾车、日常生活中的扫视都是用眼动的方式来完成的。脑卒中以后，特别是对于那些扫视潜伏期变长的人来说，他们的视线很难准确地落在某个地方。当患者的扫视变得缓慢或者容易感到疲劳的时候，就会出现对周围环境信息获取的速度和准确度下降的现象，并且和跌倒的风险以及社交活动的开展有着密切的关系。

平稳追踪是眼球追随一个不断变化的目标时所表现出来的稳定状态，即跟踪行驶中的车、飞过的鸟等物体。追踪和扫视的区别在于前者属于弹道式的运动，是由额叶眼区以及上丘来进行控制的；而后者则是需要依靠视网膜上的目标速度去进行提取，并且还需要借助小脑前庭通路发出的反馈信息来加以调节。脑卒中后平滑追踪障碍常常表现为眼球不能跟着目标的速度和方向移动，出现多次纠正性的补偿扫视。

视觉调节是眼睛从一个物体转向另一个物体时能够自动调整焦距的过程，是由睫状肌控制以及瞳孔调节共同作用来实现的。老年人由于生理老化而出现的老花眼，在脑卒中的影响下也会使患者的眼球调节功能变差。本装置用近距LED、远距激光交替发光的方式模拟日常生活里远近交替看东西的情况，从而对患者的睫状肌调节功能进行有目的性的训练。

空间注意指对视野两侧的空间信息进行均衡感知、分配的能力。单侧空间忽略是脑卒中后最常见的认知功能障碍之一，表现为患者对一侧（通常为左侧）空间的视觉和触觉信息几乎完全丧失感知，并非因视野缺损或感觉丧失，而是因为对侧空间的注意力机制受损。此项训练通过反复将视觉刺激和听觉提示施加于被忽略的视野区域，促进患者双侧空间注意力的重建，这对于偏盲和空间忽略患者尤为关键。

另外由于脑卒中病人经常会出现感觉、运动以及认知功能下降的情况，因此交互方式应当选择语音命令加简单的单键按键的形式进行操作，不能使用屏幕菜单导航或者精确的四肢协调的方式来进行任务执行。对住院和居家两种不同的使用环境设置了完整的训练模式（五个模式连续）和自选训练模式（患者自己挑选一项或多项活动练习）。

\subsection{功能模块划分}

根据以上的临床需求分析，该系统所有的功能都可以归纳成四个正交的功能层次，即感知、决策、执行和反馈。其层次结构如图\ref{fig:system_layers}所示。感知层的任务就是用惯性传感器实时采集佩戴者头部欧拉角(俯仰、横滚、侧倾)的变化情况，并且利用语音模块来捕捉到患者的口述训练指令。决策层会依照目前的姿态以及语音事件来进行系统状态机的状态转变工作，并且还会依照着训练方式及进度状况来计算出激光的目标位置以及试验的各项参数。执行层把决策的结果变成物理上的表现形式，即舵机云台控制激光点在墙面上的位置移动，语音交互集成模块发出训练引导语和反馈评价语音，蜂鸣器发出报警声。反馈层就是通过WS2812RGBLED的状态颜色、闪烁模式向医护人员或者陪同人员给出持续的行为提示，并且保存下来所有的训练过程中各项评价指标(反应时间、正确率、头稳指标、代偿次数)，最后通过UART诊断帧输出完整的运行时记录。

\begin{figure}[h!]
\centering
\includegraphics[width=1.0\textwidth]{system_layers.drawio.pdf}
\caption{系统功能模块层次划分}
\label{fig:system_layers}
\end{figure}

\section{总体架构}

\subsection{系统硬件架构}

硬件结构分为六个部分，各部分之间通过各种总线进行通信，其互联关系如图\ref{fig:hardware_arch}所示。主控子系统以DM-MC02核心板(STM32H723VGT6微控制器)为中心，用SPI、UART和GPIO三个总线来组成外围通信网络。姿态感知子系统是由BMI088六轴惯性传感器通过SPI2总线向主控提供角速度和加速度原始数据构成的；视觉刺激子系统是由SG90双轴舵机云台(TIM1PWM驱动)和650nm激光模块(PA0控制)组成的，用来完成激光点在墙面二维平面上的空间定位任务；语音交互子系统是通过CI1302AI语音模块经由UART7总线实现双向命令传输的；人机反馈子系统集成了一个患者的按键(PA2，软件消抖)、一块板载蜂鸣器(PB15，TIM12PWM驱动)和一组WS2812可编程RGBLED(PA7，SPI6模拟单总线)，给病人提供触觉、听觉和视觉三种反馈方式；电源及结构子系统则对外接直流电源进行了两级降压处理，并且设置了5V动力电源门控使能功能以及所有子系统所使用的机械固定装置和防护措施。

\begin{figure}[h!]
\centering
\includegraphics[width=0.95\textwidth]{hardware_architecture.drawio.pdf}
\caption{系统硬件架构框图}
\label{fig:hardware_arch}
\end{figure}

\subsection{系统软件架构}

系统软件采用了前后台结构(Super-loop Architecture)。后台程序每隔10ms执行一次主循环，前台用中断服务程序(ISR)来处理那些需要比较高的实时性的外设事件，比如SPI2的DMA传输完成中断、USART1和UART7的收发中断、TIM1和TIM12的PWM更新中断以及BMI088传感器的外部数据准备好中断等等都包含在里面。

训练逻辑主要由一个九状态系统的状态机构成，它可以支持完整的训练模式（Full Mode）以及自选的训练模式（Custom Mode）。完整模式下五种训练模式按照固定的顺序进行，每一项都要由患者PA2按下按钮来完成衔接工作；自选模式下，患者可以通过语音命令词逐个地挑选出想要的模式，在这个过程中，所有的模式命令词都可以在训练执行的过程中、暂停的时候以及反馈评价的时候自由地切换到其他的状态上。系统使用状态机主运行函数中的命令解析引擎来统一调度全局和状态分发语音命令。所有的硬件操作——舵机角度的设定、激光的控制、LED的控制——都是通过应用层封装的函数间接执行的，在系统暂停的时候可以正确地拦截这些操作，在恢复之后又会按照暂停前的状态进行还原。

\subsection{训练流程}

完整的单次训练流程如图\ref{fig:train_flow}所示。系统上电以后首先会进入到系统的模式选择阶段，然后患者通过按键来选择要进行的训练种类。接着外围设备初始化，所有的模块都已准备好之后，主控通过串口向语音模块发送握手协议帧，等语音模块确认已经准备好后就会播报出“初始化完成”的消息，并且开始相应的训练程序。

\begin{figure}[h!]
\centering
\includegraphics[width=1.0\textwidth]{figures/train_flow.drawio.pdf}
\caption{单次训练完整流程图}
\label{fig:train_flow}
\end{figure}


完整的模式下，系统会依次执行A、B、C、D、E这五个阶段的训练过程，在每一个训练之前都会先播报系统的训练提示专用开始提示，然后进行姿态校准，接着就是开始进行训练，并且最后会给出反馈评价结果；评价结束后就会进入到下一个模式的确认状态等待患者PA2按键确认之后才会进入下一模式并且重复循环。

自选模式下系统处于空闲状态静默等待，患者说出命令词后立即执行校准、训练和反馈，然后返回到空闲状态继续等待。在训练的过程中，病人可以随时说出另外一个训练模式的命令词来立刻转换到另外一个训练模式当中去。

系统会一直监控头部的姿态是否处在安全范围内，一旦出现异常，就会立刻停止训练，并启动一个3秒钟的自动复位定时器。病人可以随时说出口令“暂停训练”、“继续训练”、“重新开始”、“跳过这个”等等全局命令来控制整个训练过程。说出唤醒词“你好小盈”，就进入到5秒语音监听窗口，在这个窗口中可以直接捕获之后发出的所有命令，如果超过五秒没有收到任何命令的话，系统就会自动回到原来的状态去执行下一步的操作。

%\section{关键技术指标}

%根据前述功能需求和系统架构方案，本装置的详细技术指标归纳如下。

%\textbf{姿态感知指标}：姿态更新速率为100Hz（10ms主循环周期），加速度计量程为$\pm$3g、输出数据率800Hz，陀螺仪量程为$\pm$2000dps、数据带宽1000Hz/116Hz。俯仰角和横滚角的静态零偏经800样本校准后优于0.5°，在单次训练时长（15$\sim$50s）内纯陀螺仪积分漂移不超过1°。头稳指标测量范围为0°至10°，分辨率为0.01°。

%\textbf{视觉刺激指标}：激光波长为650nm（红色），光功率为5mW（Class II安全等级），激光点直径在1.5m投影距离上约为3mm至5mm。舵机云台X轴（俯仰）和Y轴（偏航）的转动范围均为0°至180°，对应墙面投影覆盖面积在1.5m投影距离上约为1.2m$\times$0.9m。舵机位置更新速率为50Hz（20ms PWM周期），静态定位分辨率为0.09°（CCR单步对应角度）。

%\textbf{语音交互指标}：语音指令识别响应时间在300ms至800ms之间。唤醒词“你好小盈”的一次唤醒成功率为约90\%，支持5m范围远场拾音。MCU向语音模块发送的TTS播报帧为16进制5字节定长协议帧（$0xAA\ 0x55\ [TYPE]\ [ID]\ 0xFB$），UART接口波特率为115200bps。系统共计支持46条独立的TTS播报词和10条可识别的命令词。

%\textbf{训练模式指标}：提供两种训练模式——完整模式自动串联五种训练（注视15s→扫视8试次→追踪8航点→聚焦5循环→忽略6试次），自选模式由患者语音命令逐一选择。训练中检测到姿态超限时触发自动暂停的保护延迟小于10ms（一个主循环周期），自动恢复所需的连续安全姿态保持时长为3s。语音监听窗口超时时间为5s。

%\textbf{物理与电源指标}：装置总质量约180g（含核心板、云台、激光模组、语音模块和外壳）；供电方式为外接直流电源通过XT30接口输入12V至24V宽电压，经两级DC-DC转换后输出5V动力电源（最大2A负载）和3.3V数字逻辑电源；整机在5V供电轨满载条件下的估算功耗约为3.5W至4.5W。
